% !TeX root = ../main.tex

\chapter{Introduction}\label{chapter:introduction}

With the continuous growth of photovoltaic (PV) systems, battery energy storage systems (BESS), and electric vehicles (EVs) in residential applications, users in traditional distribution networks are gradually evolving from passive electricity consumers into prosumers capable of electricity generation, energy storage, and flexible electricity consumption. For an individual household, PV generation can reduce electricity purchased from the grid, the stationary battery can be charged during low-price periods and discharged during high-price periods, and EV charging can also be scheduled according to user mobility requirements and electricity price signals. However, when more prosumers are connected to the same low-voltage distribution network, these flexible resources can no longer be treated only as an economic optimization problem within individual prosumers. The EV charging, battery charging and discharging, and PV utilization of multiple households can modify the total net load, affect bus voltages, line loading, and transformer loading. Therefore, building energy management should  consider economic performance from the user side, EV mobility requirements of the prosumers, and total network safety.

This thesis investigates a multi-prosumer energy management problem in a low-voltage distribution network. The system consists of 3 prosumers, each equipped with its own household loads, PV generation, a stationary battery, and EV charging demand. A simulation environment based on pandapower is developed to map the net load of the prosumers to distribution network power flow calculations, thereby evaluating distribution network safety indicators such as bus voltages, line loading, and transformer loading. The control framework in this work integrates LSTM forecasting, a MADRL actor-critic controller, EV soft or hard constraint handling, and grid safety projection, and compares their performance with baseline controllers. The main objective is to analyze how different control methods can influence the total cost, battery arbitrage, EV charging cost, departure SoC feasibility, and distribution network safety.

This thesis is organized as follows: Chapter 2 reviews the literature on building energy management and different controlling methods used today, and identifies the research gaps. Chapter 3 shows the system model of the thesis, including the models of prosumer loads, PV generation, stationary batteries, EVs, electricity prices, net load, power flow, and the optimization objective. Chapter 4 introduces the methodology, including LSTM forecasting for electricity price, the MADRL observation and action design, the reward function of MADRL, safety projection, different EV constraint handling methods, and the baseline controllers. Chapter 5 describes the experimental datasets, training settings, network configuration, and evaluation metrics for the experiments. Chapter 6 presents and analyzes the experimental results of the different controllers. Finally, Chapter 7 summarizes the main findings of this thesis and discusses outlook for future research.
